\section{Referencial Teórico}
\label{sec:theory}

\subsection{Tamanho amostral por Cochran (1977)}
\label{sec:theory:cochran}

Para uma população finita de tamanho $N$ e uma variável de interesse com
proporção populacional $p$ desconhecida, Cochran~\cite{cochran1977} demonstra
que o tamanho amostral $n_0$ necessário para estimar $p$ com erro de
amostragem $e$ e confiança $1-\alpha$ em uma população infinita é:
\begin{equation}
  n_0 \;=\; \frac{z_{1-\alpha/2}^{2} \cdot p \cdot (1-p)}{e^{2}},
  \label{eq:n0}
\end{equation}
onde $z_{1-\alpha/2}$ é o quantil correspondente da normal padrão. Para
populações finitas, aplica-se a correção de população finita (FPC):
\begin{equation}
  n \;=\; \frac{n_0}{1 + (n_0 - 1)/N}.
  \label{eq:n}
\end{equation}
Adotando o cenário conservador $p=0{,}5$ (que maximiza a variância
$p(1-p)$), $e=0{,}05$ e $1-\alpha=0{,}95$ (donde $z\approx 1{,}96$),
obtém-se $n_0 \approx 384$ e, para o ToLD-Br ($N=21{,}000$), $n=378$.

\subsection{Amostragem estratificada sem reposição}
\label{sec:theory:strat}

A população $\mathcal{U}$ é particionada em $H$ estratos disjuntos
$\mathcal{U}_1, \dots, \mathcal{U}_H$ com tamanhos $N_1, \dots, N_H$ tais que
$\sum_{h=1}^{H} N_h = N$. Em cada estrato, sorteia-se uma amostra aleatória
simples sem reposição de tamanho $n_h$, com $\sum n_h = n$. O peso amostral é
$W_h = N_h / N$.

\subsubsection{Alocação ótima de Neyman}
\label{sec:theory:neyman}

A alocação \emph{ótima} (de Neyman) minimiza a variância do estimador
estratificado da média $\bar{Y}$ para um orçamento total fixo $n$,
distribuindo $n_h$ proporcionalmente a $N_h \sigma_h$, onde $\sigma_h$ é o
desvio-padrão de $y$ no estrato $h$~\cite{cochran1977}:
\begin{equation}
  n_h \;=\; n \cdot \frac{N_h \sigma_h}{\sum_{i=1}^{H} N_i \sigma_i}.
  \label{eq:neyman}
\end{equation}
Bolfarine \& Bussab~\cite{bolfarine2005} chamam atenção para o caso especial
em que $\sigma_h$ é constante em $h$: a equação~(\ref{eq:neyman}) reduz-se à
alocação \textit{proporcional} $n_h = n \cdot W_h$.

\paragraph{Variância e ganho de Neyman.} A variância do estimador
estratificado da média sob alocação ótima é
(Cochran~\cite{cochran1977}~§5A.7; Lohr~\cite{lohr1999}~§4.5):
\begin{equation}
  \mathrm{Var}_{\mathrm{Ney}}(\hat{\bar{Y}}_{st})
  \;=\; \frac{1}{n}\Bigl(\sum_{h} W_h \sigma_h\Bigr)^{2}
        - \frac{1}{N}\sum_{h} W_h \sigma_h^{2}.
  \label{eq:var-neyman}
\end{equation}
Sob alocação proporcional, a variância correspondente é
$\mathrm{Var}_{\mathrm{prop}} = \frac{1}{n}\sum_h W_h \sigma_h^{2}
- \frac{1}{N}\sum_h W_h \sigma_h^{2}$. A diferença, isto é, o ganho que
Neyman oferece sobre a alocação proporcional, é:
\begin{equation}
  \mathrm{Var}_{\mathrm{prop}} - \mathrm{Var}_{\mathrm{Ney}}
  \;=\; \frac{1}{n}\sum_{h} W_h \bigl(\sigma_h - \bar{\sigma}\bigr)^{2},
  \qquad \bar{\sigma} = \sum_h W_h \sigma_h.
  \label{eq:gain-neyman}
\end{equation}
A equação~(\ref{eq:gain-neyman}) é o teorema central justificando a alocação
ótima: o ganho é não-negativo, é zero apenas quando os $\sigma_h$ são todos
iguais, e cresce com a heterogeneidade dos desvios-padrão entre estratos. Esta
é a hipótese empiricamente avaliada na Seção~\ref{sec:results}.

\subsection{Estimadores estratificados de $\bar{Y}$}
\label{sec:theory:estimators}

Sejam $\bar{y}_h$ e $\bar{x}_h$ as médias amostrais de $y$ e de uma variável
auxiliar $x$ no estrato $h$; $s_{yh}^2$ e $s_{xh}^2$ as variâncias amostrais;
$s_{xyh}$ a covariância amostral; $\bar{X}_h$ a média populacional de $x$ no
estrato (suposta conhecida); $\bar{X}$ a média populacional global; e
$f_h = n_h / N_h$ a fração amostrada no estrato.

\subsubsection{Média estratificada simples (\textit{baseline})}

\begin{equation}
  \hat{\bar{Y}}_{st} \;=\; \sum_{h=1}^{H} W_h\,\bar{y}_h,
  \qquad
  \widehat{\mathrm{Var}}(\hat{\bar{Y}}_{st}) \;=\;
   \sum_{h=1}^{H} W_h^{2}\,(1 - f_h)\,\frac{s_{yh}^2}{n_h}.
  \label{eq:strat-mean}
\end{equation}

\subsubsection{Estimador tipo razão (Cochran §6.10--6.11)}

A versão \emph{separada} aplica a razão dentro de cada estrato e agrega:
\begin{equation}
  \hat{\bar{Y}}_{Rs} \;=\; \sum_{h} W_h\, \frac{\bar{y}_h}{\bar{x}_h}\, \bar{X}_h.
  \label{eq:ratio-sep}
\end{equation}
A versão \emph{combinada} agrega antes da razão:
\begin{equation}
  \hat{\bar{Y}}_{Rc} \;=\; \frac{\sum_h W_h \bar{y}_h}{\sum_h W_h \bar{x}_h}\,\bar{X}.
  \label{eq:ratio-comb}
\end{equation}
A variância de grande amostra para a versão separada é~\cite{cochran1977}:
\begin{equation}
  \widehat{\mathrm{Var}}(\hat{\bar{Y}}_{Rs}) \;\approx\;
   \sum_{h} W_h^{2}\,\frac{(1 - f_h)}{n_h}\,
     \bigl(s_{yh}^{2} - 2\,R_h\,s_{xyh} + R_h^{2}\,s_{xh}^{2}\bigr),
  \label{eq:ratio-var}
\end{equation}
com $R_h = \bar{y}_h / \bar{x}_h$. Expressão análoga, com $R$ global no lugar
de $R_h$, vale para a versão combinada.

\subsubsection{Estimador tipo regressão (Cochran §7.10--7.11)}

A versão \emph{separada} usa o coeficiente de regressão dentro de cada estrato
$b_h = s_{xyh}/s_{xh}^{2}$:
\begin{equation}
  \hat{\bar{Y}}_{lrs} \;=\; \sum_{h} W_h\,\bigl[\bar{y}_h + b_h(\bar{X}_h - \bar{x}_h)\bigr].
  \label{eq:reg-sep}
\end{equation}
A versão \emph{combinada} estima um coeficiente $b$ ponderado pelos pesos
estratificados:
\begin{equation}
  \hat{\bar{Y}}_{lrc} \;=\; \bar{y}_{st} + b\,(\bar{X} - \bar{x}_{st}),
  \qquad
  b \;=\;
   \frac{\sum_h W_h^{2}\,(1-f_h)\,s_{xyh}/n_h}
        {\sum_h W_h^{2}\,(1-f_h)\,s_{xh}^{2}/n_h}.
  \label{eq:reg-comb}
\end{equation}
A variância de grande amostra para a versão separada é dada por
$\widehat{\mathrm{Var}}(\hat{\bar{Y}}_{lrs}) \approx \sum_h W_h^{2}\,(1-f_h)\,
s_{yh}^{2}(1-r_h^{2})/n_h$, onde $r_h$ é o coeficiente de correlação amostral
entre $x$ e $y$ no estrato $h$~\cite{cochran1977,lohr1999}.

\subsubsection{Eficiência relativa}

A regressão é, em geral, mais eficiente que a razão quando a relação entre $y$
e $x$ é aproximadamente linear com intercepto não-nulo; a razão é equivalente
à regressão se a relação passa pela origem. Em condições idealizadas, ambos
dominam o estimador da média estratificada simples se $\mathrm{Cor}(x,y)$ é
não-nula~\cite{cochran1977,thompson2002}.

\subsection{Detecção de toxicidade no ToLD-Br}
\label{sec:theory:toldbr}

O ToLD-Br~\cite{leite2020,leite2020_thesis} consiste em $21{,}000$
\textit{tweets} em português brasileiro, anotados por três avaliadores em
seis categorias de dano: \textit{homophobia}, \textit{obscene},
\textit{insult}, \textit{racism}, \textit{misogyny}, \textit{xenophobia}.
A construção e o protocolo de anotação são descritos em detalhe na
dissertação de Leite~\cite{leite2020_thesis}; trabalhos correlatos sobre
detecção automática de discurso de ódio em português incluem a tese de
Vargas~\cite{vargas2022_thesis}. Cada anotação é um inteiro
$v_{ij} \in \{0,1,2,3\}$ representando o número de avaliadores que sinalizou
a categoria $j$ para o \textit{tweet} $i$. A binarização operacional
consagrada na literatura é $\mathbb{1}[v_{ij} \geq 2]$ (regra de maioria), mas
o esquema ordinal preserva informação sobre o grau de concordância, útil
tanto para estratificação quanto para diagnóstico de ruído de
rotulagem~\cite{leite2020,leite2020_thesis}.

\subsection{Aplicações brasileiras de teoria amostral}
\label{sec:theory:brasil}

A literatura brasileira de amostragem e estatística aplicada
oferece referências relevantes para o desenho proposto neste artigo.
Bolfarine~\&~Bussab~\cite{bolfarine2005}, em particular, sistematizam os
estimadores estratificados, razão e regressão na tradição da Escola de
Estatística da USP. Trabalhos posteriores no Brazilian Journal of Probability
and Statistics --- por exemplo, Cordeiro~et~al.~\cite{cordeiro2010} sobre
ajustes por pós-estratificação em pesquisas amostrais com pesos --- ilustram
extensões metodológicas pertinentes a auditorias de \textit{corpora}
desbalanceados como o ToLD-Br.
